% Chapter 5: Fortune 5 Deployment Patterns at Microsecond Latency
% OLDIA Thesis — wasm4pm cognition layer

\chapter{Fortune 5 Deployment Patterns at Microsecond Latency}
\label{ch:fortune5}

\section{The Operational Economics of Microsecond Reasoning}
\label{sec:economics}

The central economic fact of the OLDIA deployment model is this: a breed execution
completing in two to fifty microseconds permits one million certified reasoning
decisions per second per CPU core. This is not an architectural aspiration; it is
a measured consequence of compiling the thirteen breed implementations to WebAssembly
and executing them against the deterministic kernel.

The implications are population-scale. A Fortune~5 pharmaceutical company running
a fifty-thousand-patient clinical trial need not sample lab results and apply
reasoning retrospectively. Every result, at the moment it enters the data stream,
can receive a certified inference with a BLAKE3 receipt attesting that the declared
protocol was followed. A Fortune~5 retailer routing millions of logistics decisions
per day need not sample exceptions for post-hoc review. Every routing decision
carries a certified soundness proof. The shift is categorical: from statistical
coverage to complete coverage.

To understand what is genuinely new here, it is necessary to situate the original
systems against which the breeds are patterned. MYCIN, the Stanford rule-based
medical consultation system described by Shortliffe and Buchanan in 1975--1976,
required minutes per consultation on the hardware of its era. STRIPS, the Stanford
Research Institute Problem Solver introduced by Fikes and Nilsson in 1971, required
minutes per planning problem. SOAR, the unified cognitive architecture described by
Laird, Newell, and Rosenbloom in 1987, operated on the order of seconds per decision
cycle. In each case the practical constraint was identical: the system could reason
correctly, but it could not reason at operational tempo. Clinical deployment of MYCIN
was studied but not adopted in practice, partly because a physician consultation
taking minutes was incompatible with the pace of ward rounds. STRIPS was embedded
in the SHAKEY robot project, where plan latency was acceptable because physical
execution was the bottleneck. SOAR was deployed in training simulations where
real-time but not microsecond response was required.

The OLDIA breeds retain the formal properties of these originals---Shortliffe-Buchanan
CF algebra, STRIPS soundness and completeness, SOAR preference algebra and subgoal
mechanism---while executing at four to six orders of magnitude faster latency. The
economic argument for enterprise deployment rests entirely on this gap. At two to
fifty microseconds, certified reasoning becomes an infrastructure service rather than
a scarce expert resource.

The subsections that follow describe eight concrete deployment patterns, each grounded
in the verified capabilities of the breed implementations as tested against 655 passing
tests (480 Rust, 175 TypeScript) on the current codebase.

\section{Pattern 1: Clinical Trial Audit Infrastructure (MYCIN Class)}
\label{sec:mycin-deployment}

\subsection*{Operational Context}

A Fortune~5 pharmaceutical company running a fifty-thousand-patient trial accumulates
lab results, adverse event reports, and dosing records at a rate that makes
per-patient expert review impossible. The regulatory obligation is not merely to
produce a correct aggregate outcome but to demonstrate that every individual
decision---every dose recommendation, every adverse event classification---followed
the declared protocol. FDA and EMA submission requirements have moved toward
computational auditability: the sponsor must be able to reconstruct the reasoning
for any flagged decision during inspection.

\subsection*{Breed Capabilities Engaged}

The MYCIN breed (\texttt{production\_rules.rs}) implements the Shortliffe-Buchanan
certainty-factor algebra in full. The three sign cases are handled correctly: when
both evidence CF and rule CF are positive, the combination formula is
$CF_1 + CF_2 \cdot (1 - CF_1)$; when both are negative,
$CF_1 + CF_2 \cdot (1 + CF_1)$; when they have opposite signs, the formula uses
the denominator $1 - \min(|CF_1|, |CF_2|)$. The strict threshold is enforced at
the boundary: a combined certainty factor of exactly 0.2 is rejected, while 0.201
is accepted. This boundary behavior is verified by a dedicated oracle test
(\texttt{breed\_adversarial.rs}, MYCIN CF strict boundary case).

The DENDRAL breed provides the complementary forbid-constraint mechanism. Where
MYCIN accumulates evidence toward a recommendation, DENDRAL enforces hard
contraindications. A candidate therapy that violates a safety-critical exclusion
criterion is eliminated regardless of its plausibility score. The forbid absoluteness
guarantee is verified by test: the forbidden candidate is eliminated even when it
carries the highest plausibility among all candidates. In a clinical context this
corresponds to a hard allergy or contraindication check that cannot be overridden
by accumulated positive evidence.

\subsection*{What Is New Versus 1976}

The original MYCIN system was consulted after the fact, by a physician who had
already formed a clinical impression. It could not operate on the entire patient
population simultaneously. The OLDIA deployment pattern inverts this: the breed
is embedded in the clinical data stream and processes every incoming lab result
as it arrives. The output is not a physician decision aid but a certified audit
record.

The BLAKE3 receipt chain is the mechanism that makes regulatory submission feasible.
Each \texttt{cognition\_run} call produces a \texttt{ContractResult} with
\texttt{output\_hash = blake3(serde\_json(BreedOutput))}, where the BreedOutput
includes the OCEL log covering the execution trace. The \texttt{run\_id} is
\texttt{blake3(breed | output\_hash)}, making the receipt a cryptographic commitment
to both the breed identity and the complete execution record. The FDA submission
package can include the full receipt chain for the trial cohort, with each receipt
independently verifiable.

The OCEL conformance layer (\texttt{src/ocel.rs}, 632 lines) enforces that the
execution trace is process-lawful before a receipt is emitted. The
\texttt{validate\_ocel\_alignment} function checks the trace against the declared
lifecycle model using a DFA phase checker; \texttt{check\_temporal\_conformance}
enforces strictly monotone logical step ordering. A non-conforming execution is
refused at the \texttt{run\_breed} choke point in \texttt{wasm.rs}. The receipt
therefore attests not merely to the output hash but to the fact that the reasoning
followed a lawful lifecycle.

\section{Pattern 2: Supply Chain Plan Certification (STRIPS Class)}
\label{sec:strips-deployment}

\subsection*{Operational Context}

A Fortune~5 retailer processes millions of logistics decisions per day: routing
exceptions, substitution approvals, carrier selections, inventory reallocation
orders. The current state of the art is largely heuristic, with exceptions routed
to human reviewers on a sampling basis. Disputed decisions are reconstructed from
system logs, a process that is time-consuming and often incomplete when state has
changed between decision and dispute.

\subsection*{Breed Capabilities Engaged}

The STRIPS breed implements classical planning with soundness and completeness
proofs. The critical property for supply chain deployment is plan soundness:
the implementation simulates plan execution and verifies that each action's
preconditions are satisfied at the point of application. This is not a static
analysis but a dynamic check against the state that would exist at execution time.
The soundness test (\texttt{strips\_gps\_level\_10.rs}, \texttt{breed\_adversarial.rs})
verifies that a plan violating preconditions at execution time is refused, not
merely flagged.

This property eliminates a class of defect that is common in operational planning
systems: a route or allocation that was valid at decision time but violates a
constraint that changed between decision and execution. The STRIPS breed cannot
recommend a plan that would fail its own preconditions when executed.

\subsection*{What Is New Versus 1971}

The original STRIPS planner operated at interactive speed on a single planning
problem. The OLDIA deployment pattern operates at microsecond latency across all
decisions simultaneously. Every order exception, every substitution, every routing
decision receives a certified plan rather than a heuristic approximation.

The receipt chain converts disputed decisions from a log reconstruction problem
to a cryptographic lookup. The \texttt{run\_id} for any decision is computable
from the breed and output hash; the receipt covers the complete plan including
the precondition state at each step. A logistics partner disputing a routing
decision can be presented with the certified plan and the state evidence that
justified it, rather than a retrospective reconstruction.

\section{Pattern 3: Regulatory Preference Enforcement (SOAR Class)}
\label{sec:soar-deployment}

\subsection*{Operational Context}

A Fortune~5 investment bank operates under MiFID~II, Dodd-Frank, and a portfolio
of national regulatory regimes. The compliance challenge is not rule lookup but
conflict resolution: two regulations may impose requirements that are jointly
unsatisfiable for a particular trade structure, and the bank must demonstrate a
consistent, auditable resolution methodology. Current practice relies on legal
opinion memos that are not computationally verifiable.

\subsection*{Breed Capabilities Engaged}

The SOAR breed (\texttt{soar.rs}, 675 lines) implements the full preference algebra:
\texttt{require}, \texttt{prohibit}, \texttt{better-than}, \texttt{worse-than},
\texttt{best}, and \texttt{indifferent} preferences. The subgoal mechanism handles
tie impasses: when two preferences are in conflict and neither dominates, the
impasse triggers a subgoal with depth cap of~2, which injects resolution preferences.
This mirrors the SOAR architecture's handling of operator tie impasses via
subgoaling and chunking.

The antisymmetry proof is the key safety guarantee: circular worse-than preferences
produce deterministic termination rather than infinite regress. This is verified
by adversarial test (\texttt{breed\_adversarial.rs}, SOAR antisymmetry case): a
preference cycle terminates with a deterministic output rather than looping.

For regulatory compliance mapping, the preference algebra maps naturally to the
regulatory structure. Absolute prohibitions (MiFID~II Article~24 suitability
requirements for retail clients) become \texttt{prohibit} preferences. Suitability
requirements become \texttt{require} preferences. Client risk profile comparisons
between instruments become \texttt{better-than} preferences. When two regulatory
requirements conflict for a given trade---a situation that arises regularly at
the intersection of MiFID~II and Dodd-Frank for cross-border transactions---the
subgoal mechanism provides a documented, auditable resolution path.

\subsection*{What Is New Versus 1987}

The original SOAR architecture operated at second-scale decision cycles. The OLDIA
SOAR breed executes at microsecond latency, making per-trade compliance certification
feasible rather than sampled. Every trade, not a sample of trades, carries a certified
preference-resolution trace in its receipt.

The receipt chain is particularly significant for regulatory submissions. An examiner
from a regulatory authority presenting a specific trade for review can be given the
\texttt{run\_id}, which maps to a receipt containing the complete preference state,
the conflict detection record, the subgoal resolution path (if any), and the final
operator selection. The derivation is not reconstructed from logs; it is the original
certified computation.

\section{Pattern 4: R\&D Institutional Memory (CBR Class)}
\label{sec:cbr-deployment}

\subsection*{Operational Context}

A Fortune~5 manufacturer accumulates engineering knowledge across decades of product
development cycles. The challenge is not storage but retrieval and adaptation: an
engineer facing a new material specification problem needs to know whether a
sufficiently similar problem was solved previously and, if so, how the past solution
must be adapted to the current parameters. Current practice relies on informal
knowledge networks and document search, both of which degrade as the organization
scales and as engineers with institutional memory retire.

\subsection*{Breed Capabilities Engaged}

The CBR breed implements the full four-R cycle. Retrieve uses Jaccard similarity
over experiment parameters: the intersection over union of parameter sets provides
a similarity score with the identity guarantee that identical query and case
parameters yield a similarity of exactly 1.0. This guarantee is verified by oracle
test (\texttt{breed\_adversarial.rs}, CBR Jaccard identity case). Reuse adapts the
past solution to the new parameter configuration. Revise validates the adapted
solution against current constraints. Retain stores the new case with a BLAKE3
identifier that provides cryptographic provenance.

The BLAKE3 retained-case identifier is the mechanism that makes institutional
memory audit-grade. Each retained case carries an identifier that is
\texttt{blake3(serde\_json(case\_data))}, making the case content verifiable
against its identifier. A case cannot be silently modified after retention; any
change produces a different identifier. Over time the case base accumulates a
cryptographically chained provenance record of every engineering decision that
was retained.

\subsection*{What Is New Versus Classical CBR}

Classical CBR systems operated at interactive speed and required manual case
curation. The OLDIA CBR breed operates at microsecond retrieval latency with
automatic BLAKE3 case provenance. A new experiment result can be checked against
the entire historical case base in microseconds, with the similarity score providing
a ranked list of past solutions and the BLAKE3 identifier providing a verifiable
pointer to each.

The four-R cycle is fully automated in the breed: a single \texttt{cognition\_run}
call with a \texttt{CBR} breed tag and a \texttt{BreedInput} containing the new
case parameters traverses retrieve, reuse, revise, and retain in sequence, emitting
an OCEL log that proves each phase occurred in the declared lifecycle order.

\section{Pattern 5: Customer Signal Fusion (Hearsay Class)}
\label{sec:hearsay-deployment}

\subsection*{Operational Context}

A Fortune~5 consumer company personalizes customer interactions across digital
and physical channels. The technical challenge is not signal collection but signal
integration: clickstream data, purchase history, support interaction history,
real-time sentiment indicators, and inventory availability must be combined into
a single ranked recommendation under time constraints that preclude sequential
processing. Current architectures typically serialize signal evaluation, introducing
latency and creating ordering artifacts in the final recommendation.

\subsection*{Breed Capabilities Engaged}

The Hearsay breed (\texttt{hearsay.rs}, 584 lines) implements the blackboard
architecture with KSAR priority scheduling. Each knowledge source has a
rating computed as \texttt{certainty * trigger\_cf}; the priority queue uses
\texttt{f32::total\_cmp} for deterministic ordering, ensuring consistent scheduling
regardless of knowledge source declaration order. Stale KSAR re-rating is implemented:
when the blackboard is updated by one knowledge source, the ratings of affected
knowledge sources are recomputed before the next scheduling decision. This means
the scheduling order adapts to new information as it arrives rather than following
a static priority assignment.

The noisy-OR confidence bounds guarantee that the combined confidence of fused
signals cannot exceed 1.0, preventing confidence inflation from signal redundancy.
The multi-level propagation is verified by test (\texttt{breed\_adversarial.rs},
\texttt{test\_multi\_level\_fusion}): signals that propagate through multiple
blackboard levels maintain correct confidence bounds at each level.

The deterministic tie-break property is operationally significant for a consumer
personalization use case: when two signals produce equal confidence for two
candidate recommendations, the tie-break is deterministic and consistent. A customer
who receives a recommendation and then immediately refreshes the page receives the
same recommendation rather than a randomly selected alternative. This consistency
is verifiable from the receipt chain.

\subsection*{What Is New Versus Classical Hearsay}

The original Hearsay-II speech understanding system operated at batch speed on
fixed-vocabulary recognition problems. The OLDIA Hearsay breed operates at
microsecond latency with cryptographically certified outputs. Every personalization
decision carries a receipt that includes the KSAR scheduling order, the blackboard
state at each update, and the final confidence scores. A customer service
representative investigating a complaint about a recommendation can retrieve the
exact signals that drove it and the order in which they were applied.

\section{Pattern 6: Quality Defect Hypothesis (DENDRAL Class)}
\label{sec:dendral-deployment}

\subsection*{Operational Context}

A Fortune~5 manufacturer running high-volume production lines relies on statistical
quality sampling to detect systematic defects. The sampling approach introduces
latency between defect onset and detection, and misses defects that affect a small
fraction of units within the sampling period. The operational goal is 100\% unit
coverage with certified defect reasoning at production line tempo.

\subsection*{Breed Capabilities Engaged}

The DENDRAL breed implements enumerate-and-test with hard constraints. The
hypothesis generation phase enumerates candidate defect explanations from sensor
readings. The test phase applies plausibility scoring. The constraint elimination
phase removes candidates that violate hard exclusion criteria, regardless of
plausibility score.

The forbid absoluteness guarantee is the key property for safety-critical quality
control. A candidate defect explanation that is excluded by a hard safety criterion
is eliminated from the hypothesis set unconditionally. An adversarial test
(\texttt{gps\_dendral\_eliza\_falsification.rs}, DENDRAL forbid absoluteness case)
verifies this: the forbidden candidate is eliminated even when it carries the
highest plausibility score among all candidates. In production quality terms, this
means a defect pattern that would trigger a safety stop cannot be suppressed by
a high-confidence plausibility assignment from other sensors.

The enumerate-and-test architecture also means that the hypothesis set is explicit
and enumerable. The OCEL log for a DENDRAL execution contains the full candidate
set, the plausibility scores, the constraint violations, and the eliminated
candidates. A quality engineer reviewing a rejected unit can inspect the complete
reasoning record rather than a summary flag.

\subsection*{What Is New Versus Classical DENDRAL}

The original DENDRAL system was applied to mass spectrometry data interpretation
and operated at interactive speed. The OLDIA DENDRAL breed operates at microsecond
latency over sensor streams, enabling 100\% unit coverage versus statistical sampling.
The receipt chain creates a continuous certified quality record: every production
unit that passes the defect hypothesis test has a receipt proving which candidate
defects were evaluated, which were eliminated by hard constraints, and which were
rejected on plausibility grounds.

\section{Pattern 7: Contract Intelligence (Prolog Class)}
\label{sec:prolog-deployment}

\subsection*{Operational Context}

A Fortune~5 company with a large procurement and legal operation processes thousands
of contracts containing complex obligation structures. Current practice applies
standard clause templates with human review of exceptions. The review bottleneck
concentrates legal expertise on the exception cases while leaving the standard-clause
population largely unanalyzed. Multi-hop obligation chains---where obligation~A
is triggered by condition~B, which is defined by clause~C, which cross-references
jurisdiction~D---are particularly difficult to analyze at scale.

\subsection*{Breed Capabilities Engaged}

The Prolog breed (\texttt{prolog8/kernel.rs}, 1143 lines) implements Robinson
flat-term unification and bounded SLD resolution with positional~\texttt{?N}
variables. Horn clauses express contract rules of the form: if clause-type~$X$
and jurisdiction~$Y$ and counterparty-property~$Z$, then obligation~$Q$. Multi-hop
derivations are supported through the standard SLD resolution mechanism: a goal
that requires intermediate facts to be derived from other rules chains through
multiple resolution steps.

The grandparent test is the canonical verification that shared-variable unification
works correctly. The rule \texttt{grandparent(?0,?2) :- parent(?0,?1), parent(?1,?2)}
with facts \texttt{parent(alice,bob)} and \texttt{parent(bob,carol)} derives
\texttt{grandparent(alice,carol)} through the shared variable \texttt{?1}. This
derivation requires that the binding of \texttt{?1} to \texttt{bob} in the first
\texttt{parent} subgoal is propagated to the second \texttt{parent} subgoal.
An implementation that passes this test has functional shared-variable unification;
an implementation without it cannot derive the conclusion regardless of the
individual fact lookups.

For contract analysis, the analogous multi-hop derivation is: \texttt{indemnity\_obligation(?party, ?loss)} may derive through
\texttt{liable(?party, ?event)}, \texttt{caused(?event, ?loss)}, and
\texttt{capped(?loss, ?amount)}, each rule sourced from a different contract
clause and possibly a different jurisdiction. The Prolog breed handles this
derivation structure natively.

\subsection*{What Is New Versus Classical Prolog}

Classical Prolog interpreters operate at interactive speed and do not produce
cryptographically certified derivation traces. The Prolog breed executes at
microsecond latency and produces a receipt whose \texttt{output\_hash} covers
the complete derivation path, not merely the final conclusion. A counterparty
disputing an obligation derivation can be presented with the receipt, which
commits to every resolution step and fact instantiation in the derivation.

The bounded SLD resolution prevents infinite loop behavior on cyclic clause
structures, which are common in real contract corpora where cross-reference
cycles arise from standard clause libraries. The depth bound is enforced at
the resolution engine level and is included in the OCEL trace.

\section{Pattern 8: Goal-Directed Planning Under Constraint (GPS Class)}
\label{sec:gps-deployment}

\subsection*{Operational Context}

A Fortune~5 logistics and operations company manages thousands of concurrent
planning instances: vehicle routing, warehouse allocation, staffing optimization,
maintenance scheduling. The planning problems are multi-objective: minimize cost,
satisfy time windows, respect capacity constraints, meet service level agreements.
Current planning systems operate at batch latency, producing plans minutes to
hours after the triggering event. This latency means plans are often stale at
execution time.

\subsection*{Breed Capabilities Engaged}

The GPS breed implements means-ends analysis against a difference reduction table.
Given a current state and a goal state, the breed identifies the largest difference,
selects an operator that reduces it, applies the operator, and recurses until the
goal is reached or the depth limit is exhausted.

The difference coverage invariant is the key property for multi-objective planning.
The GPS difference table coverage test (\texttt{gps\_dendral\_eliza\_falsification.rs})
verifies that the plan produced by the breed addresses every unmet goal fact. A
plan that satisfies the highest-priority objective while leaving a lower-priority
objective unaddressed fails this test. In operational terms, this means a logistics
plan cannot satisfy the time window constraint while silently leaving the capacity
constraint unaddressed.

The depth limit prevents runaway planning on problem instances where the difference
reduction table does not converge to a solution. This bound is enforced at the
GPS resolution level and is included in the OCEL trace, making it auditable that
planning was bounded to the declared depth.

The AutoinstinctLearning breed provides a complementary capability: monotone
distance reduction. The learning breed's oracle test (\texttt{autoinstinct\_adversarial.rs})
verifies that each plan step shows non-increasing distance to the goal. This
property can be applied to GPS plan validation: a plan that increases distance
to the goal at any step is refused.

\subsection*{What Is New Versus Classical GPS}

The original GPS system, described by Newell and Simon in 1961 and 1963, operated
at interactive speed on small problem instances. The OLDIA GPS breed executes at
microsecond latency, permitting thousands of concurrent planning instances. The
OCEL log for each planning instance records the difference sequence, the operator
selections, and the state at each step, providing a complete certified record of
the planning derivation.

\section{The Composition Pattern}
\label{sec:composition}

A property of the OLDIA architecture that is not present in any of the original
systems is uniform composability. All thirteen breeds share the
\texttt{BreedInput}/\texttt{BreedOutput} contract. The input structure is
\texttt{\{ breed: string, contract: BreedInput, options?: \{ profile? \} \}};
the output structure is
\texttt{\{ status: "ok", breed, run\_id, output\_hash, replay\_pointer, options\_profile, output \}}.
This uniformity means that a pipeline connecting multiple breeds requires no
adapter code: the output of one breed is structurally compatible with the input
of the next.

A concrete composition pattern for a pharmaceutical clinical decision support
system illustrates this:

\begin{enumerate}
\item MYCIN processes the lab result and produces a therapy recommendation with a
certainty factor. The \texttt{output\_hash} of this call becomes part of the next
call's input provenance.

\item DENDRAL receives the MYCIN recommendation and eliminates candidates with
hard contraindications. Candidates that survive constraint elimination pass to
the next stage.

\item STRIPS verifies that the treatment plan is sound: that the preconditions for
each treatment step are satisfiable in sequence.

\item CBR retrieves similar past cases and adapts the plan parameters based on
prior outcomes.

\item Hearsay fuses the structured plan recommendation with real-time patient
signals (current vital signs, current medication list, current allergy status)
from a blackboard updated by clinical monitoring systems.

\end{enumerate}

Each stage in this pipeline adds to the OCEL log. The \texttt{output\_hash} at
each stage covers the OCEL log accumulated to that point, meaning the final
receipt is a cryptographic commitment to the entire reasoning chain, not merely
the final output. An auditor reviewing a treatment decision has access to the
full derivation: what MYCIN concluded and with what certainty, which DENDRAL
constraints were applied and which candidates were eliminated, whether the STRIPS
plan was sound, which CBR cases were retrieved and how the solution was adapted,
and which Hearsay signals were fused and in what scheduling order.

This composition pattern applies equally to the supply chain, regulatory compliance,
and contract intelligence use cases described in earlier sections. The receipt
chain is not a post-hoc audit mechanism bolted onto an existing architecture; it
is the architecture. Every operation at every stage is attested, and the attestation
is cryptographically chained from the first breed invocation to the last.

\section{Summary}
\label{sec:ch5-summary}

The eight deployment patterns described in this chapter share three structural
properties that distinguish them from prior deployments of the original AI systems:

\textbf{Population coverage.} Every transaction, every patient result, every
logistics decision, every contract clause is processed, not a sample. The
microsecond latency of the breed implementations makes this economically feasible
at enterprise scale.

\textbf{Certified reasoning.} Every decision carries a BLAKE3 receipt chain
attesting that the reasoning followed the declared protocol, that the execution
trace was process-lawful (OCEL conformance enforced at the \texttt{run\_breed}
choke point), and that the output is cryptographically committed to the derivation
record. The receipt is not a summary; it covers the complete execution history
including the OCEL log.

\textbf{Uniform composability.} The \texttt{BreedInput}/\texttt{BreedOutput}
contract enables multi-breed pipelines without adapter code. Composition of
MYCIN, DENDRAL, STRIPS, CBR, and Hearsay in a single receipt-chained pipeline
is structurally trivial and produces a single chained receipt covering all five
reasoning stages.

The 655 passing tests (480 Rust, 175 TypeScript) that back these claims are not
acceptance tests for isolated components. They include breed adversarial tests
with known-input oracles, mathematical property invariants, formal property tests
for GPS, DENDRAL, and ELIZA, OCEL lifecycle conformance tests, and a
thirteen-breed bit-exact double-run determinism harness. A Fortune~5 deployment
that relies on the forbid absoluteness guarantee, the STRIPS plan soundness proof,
or the SOAR antisymmetry property can point to the specific test case that
certifies the guarantee holds and that any regression would be detected.

This is the operational thesis of the OLDIA architecture: not that AI reasoning
is possible at enterprise scale, but that certified, auditable, population-covering
AI reasoning at microsecond latency is feasible today, on current hardware, with
a verified implementation.
